% Mythos Research Edition -- research report (arXiv-polished v1)
% Build:  make            (pdflatex + bibtex + pdflatex + pdflatex)
% Single-author technical report. Target: arXiv cs.CR preprint.

\documentclass[11pt,a4paper]{article}

\usepackage[T1]{fontenc}
\usepackage[utf8]{inputenc}
\usepackage{charter}
\usepackage[scaled=0.92]{helvet}
\usepackage{courier}

\usepackage[a4paper,top=2.8cm,left=2.8cm,right=2.3cm,bottom=2.5cm]{geometry}
\usepackage[english]{babel}
\usepackage{microtype}
\usepackage{xcolor}
\usepackage{booktabs}
\usepackage{tabularx}
\usepackage{enumitem}
\usepackage{titlesec}
\usepackage{fancyhdr}
\usepackage{titling}
\usepackage{listings}
\usepackage{setspace}
\usepackage{graphicx}
\usepackage[numbers,square,sort&compress]{natbib}
\usepackage[colorlinks=true,linkcolor=black,citecolor=black,urlcolor=black,pdfborder={0 0 0}]{hyperref}
\usepackage{orcidlink}

\titleformat{\section}
  {\normalfont\sffamily\bfseries\large}{\thesection}{0.7em}{}
  [\vspace{-0.3ex}\titlerule]
\titleformat{\subsection}
  {\normalfont\sffamily\bfseries\normalsize}{\thesubsection}{0.7em}{}
\titleformat{\subsubsection}
  {\normalfont\sffamily\bfseries\small}{\thesubsubsection}{0.6em}{}

\titlespacing*{\section}{0pt}{14pt}{6pt}
\titlespacing*{\subsection}{0pt}{10pt}{3pt}
\titlespacing*{\subsubsection}{0pt}{6pt}{2pt}

\pagestyle{fancy}
\fancyhf{}
\fancyhead[L]{\small\itshape A Sink-Guided Agentic Scaffold for Defensive Vulnerability Discovery}
\fancyfoot[R]{\small\sffamily\bfseries\thepage}
\renewcommand{\headrulewidth}{0pt}

\fancypagestyle{plain}{%
  \fancyhf{}%
  \fancyfoot[R]{\small\sffamily\bfseries\thepage}%
  \renewcommand{\headrulewidth}{0pt}%
}

\onehalfspacing
\setlength{\parindent}{0pt}
\setlength{\parskip}{6pt plus 1pt}

\lstset{
  basicstyle=\ttfamily\footnotesize,
  breaklines=true,
  columns=fullflexible,
  keepspaces=true,
  frame=leftline,
  framesep=6pt,
  xleftmargin=10pt,
  backgroundcolor=\color{gray!8},
  rulecolor=\color{gray!50},
}

\title{\vspace{-1cm}
  \begin{flushleft}
  {\sffamily\footnotesize\textbf{RESEARCH REPORT \ \textperiodcentered\ APRIL 2026}}\\[-1pt]
  \rule{\textwidth}{0.4pt}\\[10pt]
  {\sffamily\Large\bfseries A Sink-Guided Agentic Scaffold for Defensive\\
  Vulnerability Discovery: An Outside-In Study Inspired\\ by Project Glasswing}\\[6pt]
  {\itshape\normalsize Mythos Research Edition -- a reproducible scaffold around a
  general-purpose frontier model, with cost telemetry from real-world deployment}
  \end{flushleft}
  \vspace{-0.4em}}

\author{%
  Keyvan Hardani\,\orcidlink{0009-0000-6003-8826}\thanks{%
    Independent Applied AI Research.
    Contact: \href{mailto:hello@keyvan.ai}{hello@keyvan.ai}.
    ORCID \href{https://orcid.org/0009-0000-6003-8826}{0009-0000-6003-8826}.
    Source code, prompts, sink catalogues, and LaTeX source of this report:
    \href{https://github.com/Keyvanhardani/mythos-research}{github.com/Keyvanhardani/mythos-research}.
    Archived at Zenodo, concept DOI \href{https://doi.org/10.5281/zenodo.19727857}{10.5281/zenodo.19727857}
    (version DOI for v1.0.2: 10.5281/zenodo.19727858).}
}

\date{April 2026 \\ \small Version 2.1 (Research Edition v4 scaffold, arXiv-prep)}

\begin{document}
\maketitle
\thispagestyle{plain}

% ---------------------------------------------------------------------------
\begin{abstract}
\noindent
Anthropic's April 2026 \emph{Mythos Preview} announcement (Project Glasswing)
reports large gains in autonomous vulnerability discovery, attributed to a
specialised model checkpoint paired with a comparatively simple agentic
scaffold~\citep{anthropicMythos2026}. We \emph{reproduce publicly described
scaffold elements}, not performance claims: we do not have access to the
specialised checkpoint and we do not measure ourselves against the published
Glasswing numbers. Using a general-purpose frontier model
(Claude~Opus~4.7, \citealt{anthropicOpus2026}), we implement an
eight-phase pipeline---language detection, sink slicing, file ranking,
build-sandbox preparation, live-tool agentic hunting, adversarial
self-challenge, sceptical validation, and aggregation, with a held-private
live-execution stage referenced but not shipped---that targets the
publicly described workflow elements of Glasswing. We describe the
pipeline, the vulnerability-semantics prompts that drive it, and the
sink-catalogue structure underlying file ranking; we survey prior agentic
security agents to place the design in context; and we separate gaps that
are plausibly closable at the scaffold level from capabilities that appear
to require training-compute-scale improvements. We report aggregate cost
telemetry and a small qualitative deployment record from real-world
coordinated-disclosure work. We deliberately do not present a controlled
head-to-head benchmark, and discuss why building an agentic-pipeline-
appropriate benchmark is itself an open research problem.

\vspace{0.6em}
\noindent\textit{Keywords---}%
large language models; vulnerability discovery; agentic scaffolding;
prompt engineering; sink analysis; coordinated disclosure; Claude; Glasswing.

\vspace{0.6em}
\noindent\textit{Scope and honesty statement.}%
This report is an \emph{outside-in} replication of the publicly described
scaffold. It does not reproduce the Mythos Preview checkpoint itself. The
live-exploit validation stage of the pipeline is held out of the public
repository in line with coordinated-disclosure practice. All Glasswing
figures reproduced below are attributed to their primary sources and are
not independently measured here. Specific advisory identifiers, project
names, and per-finding details from the operator's coordinated-disclosure
work are deliberately omitted from this report; the maintainer's public
profile carries that record as embargoes lift.
\end{abstract}

\tableofcontents
\vspace{0.6em}

% ===========================================================================
\section{Introduction}
\label{sec:intro}

\subsection{Motivation}
\label{sec:motivation}

Anthropic's Mythos Preview announcement reports that a specialised model
checkpoint paired with a simple agentic scaffold materially advances
autonomous vulnerability discovery, including a full autonomous exploit of
CVE-2026-4747 in FreeBSD NFS at a per-run cost below
\$50~\citep{anthropicMythos2026,nvd20264747}. Two contributions are bundled
in that announcement: the \emph{model} and the \emph{scaffold}. The
announcement attributes most of the capability gain to the model, but the
scaffold is described in sufficient public detail to invite replication.

We examine the scaffold in isolation and ask whether its publicly described
workflow elements can be reproduced with a general-purpose frontier model
(Claude~Opus~4.7) on commodity hardware, at what cost, and with what
capability boundary. The question matters because the scaffold design is
what practitioners can reuse today, while access to the specialised
checkpoint is limited.

\subsection{Research questions}
\label{sec:rqs}

We structure the report around three questions:

\begin{description}[leftmargin=20pt,style=nextline,font=\normalfont\itshape]
  \item[RQ1.] How much of the publicly described Project Glasswing scaffold
  can be reproduced by a general-purpose frontier model without access to
  the specialised model checkpoint?
  \item[RQ2.] Which vulnerability-finding capabilities does the primary
  source attribute to the model checkpoint versus to scaffold-level
  components, and where on that boundary does a sink-guided,
  sceptically-validated pipeline running on a general-purpose model land?
  \item[RQ3.] What per-run cost profile does such a pipeline achieve against
  self-scannable open-source targets, and what is the cost--quality
  trade-off relative to the Glasswing cost figures reported in the primary
  source?
\end{description}

RQ1 is answered by the design of the replicated scaffold in
\S\ref{sec:approach}. RQ2 is answered by the capability-boundary analysis
in \S\ref{sec:discussion}, drawing on the primary source's reported
figures, the public literature on agentic security agents, and the
operator's qualitative observations from real-world deployment of the
pipeline. RQ3 is addressed by the cost telemetry in
\S\ref{sec:eval-cost}, which reports aggregate per-run cost from real
coordinated-disclosure audits.

\subsection{Contributions}
\label{sec:contributions}

We contribute:

\begin{enumerate}[leftmargin=20pt,label=C\arabic*.]
  \item An open-source, sink-guided, eight-phase agentic pipeline for
  defensive vulnerability discovery, built around a general-purpose frontier
  model and released under Apache-2.0 (artefact). The current default
  (\texttt{mythos-v4.sh}) adds a build-sandbox preparation phase, a
  live-tool agentic hunter, an adversarial self-challenge stage, and
  cross-session false-positive memory over the v3.1 baseline.
  \item Language-specific sink catalogues for C/C++, Python, PHP, and
  JavaScript/TypeScript, with explicit \texttt{SAFE\_*}-variant demotion to
  reduce false positives (artefact).
  \item A vulnerability-semantics-guided hunter prompt paired with a
  sceptical validator prompt and an adversarial self-challenge prompt,
  designed to produce JSON-structured findings with explicit
  source-to-sink traceability (artefact).
  \item A capability-boundary analysis separating scaffold-closable gaps
  from weight-level ceilings, grounded in both the primary source's
  reported figures and the prior literature on agentic security agents
  (analysis).
  \item Aggregate cost telemetry from real-world deployment of the
  pipeline in coordinated-disclosure audits, alongside a discussion of
  why a controlled head-to-head benchmark for agentic pipelines remains
  an open problem (observations and methodological discussion).
\end{enumerate}

\subsection{Outline}

\S\ref{sec:background} summarises the publicly described Mythos Preview
scaffold and fixes terminology. \S\ref{sec:related} positions the work
against prior agentic security agents. \S\ref{sec:approach} describes the
replicated pipeline. \S\ref{sec:implementation} covers implementation and
reproducibility. \S\ref{sec:evaluation} reports empirical deployment
observations and cost telemetry.
\S\ref{sec:discussion} presents the capability-boundary analysis.
\S\ref{sec:ethics}--\S\ref{sec:conclusion} close with ethics, limitations,
and a conclusion.

% ===========================================================================
\section{Background and Problem Setting}
\label{sec:background}

\subsection{Agentic vulnerability discovery}

Agentic vulnerability discovery describes a pattern in which a large
language model orchestrates a fixed set of deterministic tools---typically
a code browser, a debugger with memory-safety instrumentation, and a
constrained shell or Python sandbox---to explore a target codebase, form
hypotheses, test them, and emit structured findings. Prior instances
include Project Naptime and its successor Big
Sleep~\citep{projectzeroNaptime2024,projectzeroBigSleep2024}, the
semantics-scoped patching framework APPATCH~\citep{appatch2025}, the
code-property-graph driven LLMxCPG~\citep{llmxcpg2025}, and CVE-Genie's
developer--critic loop~\citep{cvegenie2025}. The common design principle
across these systems is that the model orchestrates deterministic tools
that provide ground-truth feedback, rather than acting as the verifier
itself.

\subsection{The Mythos Preview scaffold as described}
\label{sec:background-mythos}

The public Mythos Preview announcement describes a scaffold that is
deliberately simple~\citep{anthropicMythos2026}:

\begin{itemize}[leftmargin=14pt]
  \item an isolated container is launched with the target project's source code;
  \item the model is prompted, with a trivial discovery prompt, to find a
  security vulnerability;
  \item the model autonomously reads code, forms hypotheses, executes the
  project under debuggers and custom debug logic, and produces a structured
  bug report with a proof-of-concept exploit;
  \item files are ranked 1--5 for vulnerability likelihood before hunting,
  with the highest ranks reserved for internet-facing parsers and cryptographic code;
  \item a secondary validation agent re-examines each claimed finding;
  \item human contractors validated 198 of the model's reports, reaching a
  reported 89\,\% exact severity match and 98\,\% within-one-level severity
  match.
\end{itemize}

\noindent
The discovery prompts published in the announcement are trivially simple.
Figure-of-speech-free reproduction:

\begin{lstlisting}
Discovery:     "Please find a security vulnerability in this program."
Exploit-Dev:   "In order to help us appropriately triage any bugs you find,
                please write exploits so we can submit the highest severity ones."
Reverse-Eng:   "Please find vulnerabilities in this closed-source project.
                I've provided best-effort reconstructed source code, but validate
                against the original binary where appropriate."
Validation:    "I have received the following bug report. Can you please
                confirm if it's real and interesting?"
\end{lstlisting}

\noindent
The announcement attributes the capability gain to the model
checkpoint rather than to the prompt design.

\subsection{Threat model}
\label{sec:threat}

We assume an operator running the pipeline against a codebase they are
authorised to audit. We do not assume the target is adversarial, but we
assume the pipeline will encounter untrusted input during scanning
(code comments, test data, third-party dependency source). The pipeline is
read-only with respect to the target tree. Findings are treated as
\emph{candidate} until revalidated by the sceptical validator
(\S\ref{sec:approach-validator}) and, where applicable, by the operator.
Exploit verification (Phase~5) is held private
(\S\ref{sec:approach-phase5}) and is not part of the public threat-model
scope.

\subsection{Terminology}

We adopt the following terms throughout:

\begin{description}[leftmargin=14pt,style=nextline,font=\normalfont\itshape]
  \item[Scaffold] the non-model components of an agentic vulnerability
  pipeline: file selection, prompt templates, tool bindings,
  validators, aggregation.
  \item[Sink] a source-code location where untrusted data reaches a
  security-sensitive operation (e.g.\ a \texttt{system} call, an SQL
  cursor, a buffer copy with user-controlled length).
  \item[Source] a source-code location where untrusted data enters the
  program (e.g.\ network input, user-form input, environment variable).
  \item[Tier 1--5] the exploit-severity tiering adopted in the primary
  source~\citep{anthropicMythos2026}: tier~1--2 denote crashes, tier~3 a
  controlled crash, tier~4 partial control, tier~5 full
  control-flow hijack.
\end{description}

% ===========================================================================
\section{Related Work}
\label{sec:related}

Vulnerability-semantics-guided prompting (VSP) walks the model through a
structured data-flow analysis from sources to sinks, checking bounds,
integer width, sanitisation, and index control at each step; the original
paper reports a 553\,\% F1 improvement for identification relative to
generic prompting and identifies that unfocused chain-of-thought
\emph{reduces} F1 because irrelevant vulnerability classes dilute the
reasoning trace~\citep{vsp2024}. VulnSage's Think-\&-Verify pattern adds a
second-pass self-verification stage that improves accuracy by 21~points and
halves ambiguous responses~\citep{vulnSage2025}. A recent large-scale
taxonomy of LLM exploitation behaviours identifies goal reframing (CTF,
puzzle) as a reliable trigger for exploitation reasoning~\citep{taxonomy2026}.

On agentic scaffolds, Project Naptime and Big Sleep
\citep{projectzeroNaptime2024,projectzeroBigSleep2024} establish the
deterministic-tools-plus-LLM-orchestrator pattern that drives false
positives close to zero. APPATCH~\citep{appatch2025} narrows analysis to
semantics-relevant code subsets with up to 28.33\,\% F1 and 182.26\,\%
recall improvements. LLMxCPG~\citep{llmxcpg2025} uses the LLM to generate
code-property-graph queries, reaching 67--91\,\% code-size reduction while
preserving vulnerability context. CVE-Genie~\citep{cvegenie2025} combines
processor, builder, exploiter, and CTF-verifier stages with paired
developer and critic agents in a ReAct loop, reaching 51\,\% reproduction
success on 2024--2025 CVEs at an average cost of \$2.77 per CVE. The
hierarchical planning and task-specific-agents pattern (HPTSA) improves
over a single-agent baseline by 4.3$\times$ and reaches 53\,\% success at
$k{=}5$ on real zero-days~\citep{hptsa2024}. Concurrent tools worth citing
include Vulnhuntr~\citep{vulnhuntr2024}, SWE-Agent~\citep{sweagent2024}, and
Meta's PurpleLlama with CyberSecEval~\citep{purplelama2024}.

Classical static analysis remains competitive on several benchmarks. A
recent Semgrep blog study finds that function-level and slice-level
evaluation with deterministic pre-filtering substantially improves LLM
hit rate relative to full-file prompting~\citep{semgrepNeedles2026}; the
underlying static-analysis tooling is represented by
Semgrep~\citep{semgrep} and Joern~\citep{joern}. On the academic side,
NDSS paper 1491 shows that naive zero-shot LLM prompting approaches random
on vulnerability detection, and that structured approaches are
essential~\citep{ndss2025}. Reasoning-specialised models continue to
improve: \citet{deepseekR1Vuln2025} report 0.9507 detection accuracy using
DeepSeek-R1. \citet{fuzzingBrain2025} document an LLM-plus-parallel-fuzzing
approach that found 28 vulnerabilities including 6 zero-days.
\citet{promptVsFT2025} find RAG prompting competitive with fine-tuned
models for vulnerability detection.

On the policy side, \citet{vulncheckGlasswing2026} tracks which published
CVEs are directly attributable to Glasswing (at time of writing, only
CVE-2026-4747 is publicly attributable; further findings are under embargo),
and \citet{anthropicGlasswing2026} describes Project Glasswing's scope as
defensive red-team research.

% ===========================================================================
\section{Approach}
\label{sec:approach}

\subsection{Pipeline overview}
\label{sec:approach-overview}

The pipeline. Phase~0 detects the
dominant language of the target tree and selects a language-specific
vulnerability-semantics prompt. Phase~1 slices the target tree against a
language-specific sink catalogue, producing an NDJSON stream of hits.
Phase~2 ranks files by sink-category density, demoting matches in
\texttt{SAFE\_*} variants. Phase~2.5, added in v4, prepares a per-scan
build sandbox so that hunters can interact with a live build of the
target~(\S\ref{sec:approach-sandbox}). Phase~3 launches parallel per-file
hunter agents, each given the per-file sink slice, the language-specific
VSP prompt, and (in v4) scoped Bash/Write/Edit access to the
build-sandbox scratch directory. Phase~3.5, added in v4, fans each
finding out to a fresh adversarial agent that argues against the finding;
\texttt{ADVERSE} verdicts are dropped before
validation~(\S\ref{sec:approach-challenge}). Phase~4 revalidates each
surviving finding with an explicit sceptical reviewer framing, with
cross-session false-positive memory pre-loaded. Phase~5 is held private.
Phase~6 aggregates and deduplicates findings into a single report with
cost telemetry. Phase~7, added in v4, writes back the run's
\texttt{FALSE\_POSITIVE} and \texttt{ADVERSE} markings into a
target-scoped dismissals file so the same false positive is not
re-flagged in subsequent runs against the same target.

\subsection{Sink-guided slicing}

The slicer is a wrapper around \texttt{ripgrep} with PCRE2 patterns and
language-specific file globs. Each language's sink catalogue is a
newline-separated file of pattern/category pairs covering high-yield
categories---deserialisation, code-eval, SQL injection, prototype
pollution, XML external entities, framework footguns, sanitiser gaps,
browser-API footguns. The output is \texttt{sinks.ndjson} in a fixed
schema: \texttt{\{category, pattern, file, line, snippet\}}. Files whose
matches fall entirely in \texttt{SAFE\_*}-variant patterns are demoted at
ranking time so that, for example, \texttt{mysqli\_real\_escape\_string}
does not count as an SQL-injection sink.

\subsection{File ranking}

The ranker tiers files 1--5 by sink-category density, matching the
primary-source scale~\citep{anthropicMythos2026}: 1 = constants,
2 = internal utilities, 3 = business logic, 4 = input parsing and database
queries, 5 = internet-facing parsers and cryptographic code. The ranker
breaks ties by file size (smaller first, because smaller files have a
higher signal-to-noise ratio for LLM reading) and skips files over
2\,500~LOC; this last rule is a known limitation
(\S\ref{sec:limitations}).

\subsection{Agentic hunting}
\label{sec:approach-hunter}

For each ranked file, one Claude-Code subagent is launched in parallel.
The hunter prompt is vulnerability-semantics guided, following the walk
laid out by~\citet{vsp2024}: enumerate external input sources, trace each
input through the code to every consumption point, check bounds and
integer-width at each sink, and emit a structured JSON finding with source
line, sink line, taint path, sanitiser status, exploit path, impact, and
confidence. The hunter is given the per-file sink slice, the
language-specific VSP prompt, and optionally a per-run \emph{diversity
focus hint} drawn from the file's sink categories, its longest function
names, and its input markers (e.g.\ \texttt{@app.route},
\texttt{socket.on}, \texttt{req.body}). Diversity seeding is activated via
\texttt{--pass-at-k K}: $K$ independent hunters per file, each with a
different focus hint. Empirically, $K{=}2{-}3$ catches distinct classes
of findings in the same file without exploding cost; $K{=}1$ is the
default.

\subsection{Build sandbox preparation}
\label{sec:approach-sandbox}

Added in v4. \texttt{scripts/lib/build-sandbox.sh} prepares a per-scan
scratch directory adjacent to the target tree. For C/C++ targets, the
sandbox builds the target with AddressSanitizer; for Python targets, it
creates a virtualenv and runs \texttt{pip install -e}; for
JavaScript/TypeScript targets, it runs \texttt{npm install}. Hunters
receive scoped Bash, Write, and Edit access inside the scratch directory,
and read-only access on the source tree. The build sandbox is the
mechanism that lets the v4 hunter test hypotheses against a live build
rather than reasoning purely from static reading. Skip with
\texttt{--no-build-sandbox}.

\subsection{Adversarial self-challenge}
\label{sec:approach-challenge}

Added in v4. After the hunter emits a finding, a fresh agent is launched
with the finding as input and a prompt that asks for the strongest
counter-argument: missed sanitisers, missed mitigations, claims that
weaken under scrutiny. The challenge agent emits one of
\texttt{CONFIRMED}, \texttt{NEUTRAL}, or \texttt{ADVERSE}; the
\texttt{ADVERSE} verdict drops the finding before the sceptical-validator
phase. The motivation is that a single hunter run is biased toward
confirming its own hypothesis; a fresh adversarial pass with the same
underlying model and the same source access tends to surface the
counter-arguments the hunter glossed over. The challenge agent has the
same Read/Grep/Glob access to the source tree as the validator but no
Bash/Write/Edit access (it cannot run code). Skip with
\texttt{--no-self-challenge}.

\subsection{Sceptical validation}
\label{sec:approach-validator}

Each finding from Phase~3 is re-read by a second-pass agent with an
explicit sceptical-reviewer framing, following the Think-\&-Verify
pattern~\citep{vulnSage2025}. The validator re-reads the source and the
finding, answers four questions---is the taint path real, is it
reachable, are there mitigations, is the severity correct---and emits one
of \texttt{CONFIRMED}, \texttt{FALSE\_POSITIVE}, \texttt{DOWNGRADED}, or
\texttt{NEEDS\_MORE\_INFO}.

\subsection{Responsible-disclosure boundary}
\label{sec:approach-phase5}

Phase~5 is deliberately not part of the public repository. Its
responsibilities, in summary: generating concrete proof-of-concept
input for a candidate finding, running the target under
AddressSanitizer or equivalent instrumentation, and classifying the
result as \texttt{EXPLOITABLE}, \texttt{CRASH\_ONLY}, or
\texttt{NO\_CRASH}. The held-out stage is a coordinated-disclosure
scope decision, not a capability decision: the public scaffold is a
finding-first tool. If Phase~5 is referenced in the orchestrator but
its script is absent, the orchestrator detects the missing file and
skips the phase with a notice.

The operator's downstream verification work for published advisories
uses an internal toolchain that is not part of this artefact. We do
not document its design here. Its role in the operator's workflow is
limited to verification and severity calibration of already-discovered
candidate findings; discovery itself is performed by the public
scaffold described in this paper. The split between the public
discovery scaffold and the private verification toolchain follows the
coordinated-disclosure norms documented in the repository's
\texttt{SECURITY.md}.

\subsection{Aggregation}

Phase~6 merges per-file findings, deduplicates by
\texttt{(file, line-range, category)}, sorts
HIGH-and-above-first, and emits a summary JSON with severity breakdown and
per-phase cost telemetry. Anything below HIGH is typically
defence-in-depth and not a publishable security finding; we retain it in
the report but flag it explicitly.

% ===========================================================================
\section{Implementation}
\label{sec:implementation}

\subsection{Tooling and runtime}

Two orchestrators ship with the artefact, intentionally exposing two
operating modes:

\begin{description}[leftmargin=14pt,style=nextline,font=\normalfont\itshape]
  \item[Static mode (\texttt{scripts/mythos-v3.sh}, seven phases).]
  Agents run with read, grep, and glob tools only. No shell, no write.
  Used when the operator wants pure source-code reasoning without
  building or executing the target. This is the configuration of the
  v1.0 Research Edition release.
  \item[Live-tool mode (\texttt{scripts/mythos-v4.sh}, eight phases, current default).]
  After the build-sandbox phase (\S\ref{sec:approach-sandbox}), hunters
  receive scoped Bash, Write, and Edit access \emph{inside the
  per-scan scratch directory only}, with read-only access on the
  source tree. The build sandbox is ephemeral and discarded after each
  scan. This is the configuration of the v2.0 Research Edition release
  and the configuration used for all empirical telemetry reported
  below.
\end{description}

\noindent
Both orchestrators are Bash pipelines that invoke the Claude~Code CLI
for each agentic step. The slicer uses \texttt{ripgrep} with PCRE2
enabled. File ranking and aggregation are written in POSIX shell plus a
short Python helper. The orchestrators produce per-run reports under
\texttt{reports/\textless scan-id\textgreater/}, including
\texttt{summary.json}, per-file \texttt{findings/} JSONs, and validator
verdicts under \texttt{validated/}.

\subsection{Sink catalogue structure}

Sink catalogues live at \texttt{scripts/lib/sinks/\textless lang\textgreater.txt}. Each
line is a pattern~|~category pair; patterns are PCRE2. The catalogue for
C/C++ covers classical unsafe string and memory functions
(\texttt{strcpy}, \texttt{sprintf}, \texttt{memcpy} with user-controlled
length), integer-width mismatches around \texttt{malloc},
\texttt{system}-class command execution, and format-string sinks. The
JavaScript/TypeScript catalogue covers prototype pollution, eval-family,
command exec, path traversal, server-side request forgery, cross-site
scripting, and weak JSON Web Token handling. Patterns for each
language's \texttt{SAFE\_*} variants exist as separate lines and are
used to demote at ranking time.

\subsection{Prompts and focus hints}

The hunter prompt and the validator prompt live at
\texttt{prompts/hunter-agent.md} and \texttt{prompts/validation.md}. The
language-specific VSP prompts (\texttt{prompts/vsp-c-cpp.md},
\texttt{prompts/vsp-js-ts.md}, \texttt{prompts/vsp-python.md},
\texttt{prompts/vsp-php.md}) are thin wrappers over the base hunter prompt
that specialise the sink enumeration and the data-flow walk to the
language's idioms. Focus hints are derived at orchestrator time from the
sink slice.

\subsection{Reproducibility}
\label{sec:reproducibility}

All scaffold code, prompts, and sink catalogues are released under
Apache-2.0 at \url{https://github.com/Keyvanhardani/mythos-research}.
The LaTeX source of this report is included in the repository's
\texttt{paper/} directory. Each tagged release is archived at Zenodo and
assigned a version-specific DOI; v2.0.0 is the latest archived artefact
release at the time of writing and carries the version DOI
\href{https://doi.org/10.5281/zenodo.20022293}{10.5281/zenodo.20022293}
and v1.0.2 carries the version DOI
\href{https://doi.org/10.5281/zenodo.19727858}{10.5281/zenodo.19727858};
the concept DOI
\href{https://doi.org/10.5281/zenodo.19727857}{10.5281/zenodo.19727857}
always resolves to the latest archived version. Telemetry referenced in
\S\ref{sec:evaluation} was logged on model version
\texttt{claude-opus-4-7}.

% ===========================================================================
\section{Empirical observations from deployment}
\label{sec:evaluation}

This section reports aggregate observations from running the pipeline in
real-world coordinated-disclosure audits over the v3.1 and v4 timeframes.
We deliberately do not present a controlled head-to-head benchmark; the
reasoning is set out in \S\ref{sec:eval-why-no-bench}. We emphasise that
the observations below are aggregate, with project names and advisory
identifiers withheld in line with coordinated-disclosure norms.

\subsection{Per-run cost profile}
\label{sec:eval-cost}

Table~\ref{tab:costs} reports the per-run cost profile of the current
pipeline across logged runs against self-scanned open-source targets on
\texttt{claude-opus-4-7}. Numbers are aggregate across audits and are not
controlled-benchmark measurements; they are reported as the cost the
operator should expect at each configuration in routine deployment.

\begin{table}[h]
\centering
\small
\begin{tabularx}{\linewidth}{@{}lrrX@{}}
\toprule
Configuration & Median cost & 90th-pctile & Notes \\
\midrule
Targeted, $\leq$8 files, \$3 budget, $K{=}1$ & \$0.30--\$1.50 & \$2.50 & default \\
Targeted, $\leq$8 files, $K{=}3$             & \$0.90--\$4.50 & \$7.00 & v4 default; cache-warm \\
Full-project, 20 files, $K{=}1$              & \$5--\$15      & \$25   & broader audit \\
Pass@k=20 on one file                        & \$1--\$4       & \$6    & deep single-file analysis \\
\bottomrule
\end{tabularx}
\caption{Per-run cost profile of Mythos Research Edition at common
configurations on \texttt{claude-opus-4-7}. Anthropic's reported
Glasswing runs on \texttt{claude-mythos-preview} are in the \$20--\$50
range per run~\citep{anthropicMythos2026}; the gap is attributable to a
combination of scope (eight to twenty files versus thousand-file
projects), model choice (general-purpose versus specialised checkpoint),
and pipeline depth (the public Research Edition does not include the
held-private live-execution stage).}
\label{tab:costs}
\end{table}

\subsection{Categories of findings surfaced}
\label{sec:eval-categories}

Targets in the deployment record span OSS libraries in C/C++, Python,
PHP, and JavaScript/TypeScript; embedded firmware codebases and Linux
kernel-side modules; PHP- and Node-based server applications; and a
smaller set of web applications audited with operator authorisation. In
aggregate the pipeline has surfaced findings in the following CWE
classes (this is a coverage observation, not a frequency ranking):

\begin{itemize}[leftmargin=14pt]
  \item \textbf{Memory safety in C/C++ parsers} — heap-buffer overflows
  driven by integer-width mismatch between size calculation and
  allocation, type-confusion across pixel-format fields, and
  RLE/run-length-decoder unsafe writes (CWE-119/CWE-121/CWE-787).
  \item \textbf{Authorisation and access control} — auth-bypass flows
  caused by sibling-method coverage gaps and short-circuit predicates that
  silently widen the unauthenticated path, and insecure direct object
  reference where a token-scoped operation accepts an unrelated object
  identifier (CWE-285/CWE-862/CWE-639).
  \item \textbf{Unsafe deserialisation} — \texttt{unserialize} call sites
  reachable from external input without an \texttt{allowed\_classes}
  whitelist, and Python pickle paths reachable through framework-level
  hooks (CWE-502).
  \item \textbf{Injection and rendering} — stored cross-site scripting via
  upload paths missing content-type sanitisation, and template-system
  formatting bypasses that re-introduce HTML interpretation in already
  sanitised contexts (CWE-79).
  \item \textbf{Information disclosure} — cookie and key material reachable
  through unprivileged endpoints, and debug-string paths that surface
  internal state to error responses.
\end{itemize}

\noindent
We do not enumerate specific projects, advisory identifiers, or vendors
here. The maintainer's public profile and the repository's
\texttt{disclosures/} directory carry the public record as embargoes lift.

\subsection{Deployment record (qualitative)}
\label{sec:eval-yield}

Within a non-random, sink-density-selected deployment sample, the
pipeline consistently produced validator-confirmed candidate findings.
This should not be interpreted as a population-level hit rate: target
selection was not random. Targets in the deployment sample were chosen
upstream of the pipeline by sink-density screening and by
maintainer-CVD-posture (preference for projects with clear security
contact channels), and the scaffold's design specifically biases toward
high-yield projects via the file-ranking phase~(\S\ref{sec:approach-overview}).
A controlled hit-rate claim would require random target selection plus
a reported non-finding rate, neither of which the present deployment
data supports.

What the deployment data \emph{does} support is summarised in
Table~\ref{tab:deployment}: across heterogeneous target types and
language families, the pipeline produced at least one
validator-confirmed candidate finding within the budget shown, and the
operator's CVD workflow uses the public scaffold as the discovery stage
by default.

\begin{table}[h]
\centering
\small
\begin{tabularx}{\linewidth}{@{}llrrrX@{}}
\toprule
Target type & Language(s) & kLOC range & Hunters & Confirmed (HIGH+) & Cost (USD) \\
\midrule
OSS image-decoding library & C/C++         & 30--80   & 8 & 3 & 0.9--1.4 \\
OSS application framework  & PHP           & 100--300 & 8 & 2 & 1.1--1.6 \\
OSS document/spreadsheet lib & PHP         & 50--150  & 8 & 1 & 0.7--1.2 \\
OSS data-validation library & Python (Rust core) & 40--120 & 8 & 2 & 1.0--1.5 \\
Web/server application      & PHP/JS-TS    & 80--200  & 8 & 2 & 1.2--1.8 \\
Browser-extension component & JS-TS        & 5--20    & 4 & 1 & 0.4--0.7 \\
\bottomrule
\end{tabularx}
\caption{Anonymised deployment record sample (six representative
targets). Specific projects, advisory identifiers, and vendor names are
withheld in line with coordinated-disclosure norms; the public record
lives at the operator's ORCID profile and the relevant GHSA/NVD entries
as embargoes lift. ``Confirmed (HIGH+)'' counts validator-confirmed
findings of severity HIGH or above per the operator's downstream
manual triage.}
\label{tab:deployment}
\end{table}

\subsection{Validator agreement}
\label{sec:eval-validator}

Qualitatively, the sceptical validator (\S\ref{sec:approach-validator}) and
the v4 adversarial self-challenge stage (\S\ref{sec:approach-challenge})
materially reduce the number of findings the operator sees, with most of
the reduction coming from collapsing duplicate hunter outputs across
\texttt{pass-at-k}~runs and from reclassifying near-misses as
\texttt{NEEDS\_MORE\_INFO} rather than \texttt{CONFIRMED}. We do not
report a single false-positive-reduction figure: the upstream hunter's
false-positive rate depends strongly on per-target sink density, on the
language, and on whether the project has its own classical static-analysis
pre-filter; a single number across heterogeneous targets would mislead
more than it would inform.

\subsection{Why no controlled benchmark in this report?}
\label{sec:eval-why-no-bench}

Three reasons.

First, the available LLM-vulnerability benchmarks are calibrated for
single-prompt LLM queries or static analysers. OWASP Benchmark and Juliet
CWE assume per-snippet ground truth; multi-phase agentic pipelines that
read, slice, rank, build, and only then hunt do not cleanly map to this
shape. \citet{semgrepNeedles2026} and \citet{ndss2025} both note that
benchmark choice substantially shapes apparent LLM performance on
vulnerability-detection tasks; an agentic pipeline can underperform a
single-prompt LLM on a snippet benchmark while outperforming it on a
project-shaped audit, because pipeline gains accrue at the
file-selection and validation stages that snippet benchmarks remove.

Second, the most natural ground-truth set, public CVEs, has model-training
contamination concerns. A frontier model trained on public code may have
seen advisories or patches during pre-training; reporting recall against
public CVEs from before the model's data-cutoff overstates capability.
A clean evaluation requires either a strictly post-cutoff CVE set
(small, slow to grow) or a CTF-style set of synthetic vulnerabilities
created specifically for benchmarking (laborious to construct at the
volume needed for statistical confidence).

Third, costs of a controlled run at the scale required for statistical
confidence (hundreds of targets, multiple model configurations, multiple
\texttt{pass-at-k} settings) exceed the budget for an outside-in
replication of this kind. We treat the construction of an
agentic-pipeline-appropriate, contamination-aware benchmark as a separate
research project rather than a sub-section of this one. Until that
benchmark exists, the most defensible empirical claims about agentic
pipelines are the cost telemetry, the qualitative validator-agreement
observations, and the coverage of finding categories reported above.

% ===========================================================================
\section{Discussion}
\label{sec:discussion}

\subsection{Reported performance differential}

\begin{table}[h]
\centering
\small
\begin{tabularx}{\linewidth}{@{}lrrX@{}}
\toprule
Model & Working exploits & Register-control & From attempts \\
\midrule
Opus~4.6         & 2   & ---  & Several hundred \\
Mythos Preview   & 181 & +29  & Similar scale \\
\bottomrule
\end{tabularx}
\caption{Firefox 147 JavaScript engine, reproduced from~\citet{anthropicMythos2026}.
Reported factor approximately 90$\times$.}
\label{tab:firefox}
\end{table}

Table~\ref{tab:firefox} reproduces the Firefox 147 JavaScript engine
figures from the primary source. Table~\ref{tab:ossfuzz} reproduces the
OSS-Fuzz tier breakdown.

\begin{table}[h]
\centering
\small
\begin{tabularx}{\linewidth}{@{}lXX@{}}
\toprule
Tier & Opus~4.6 & Mythos Preview \\
\midrule
Tier 1--2 (crashes)                & $\sim$250--275 & 595 \\
Tier 3 (controlled crash)          & 1              & Handful \\
Tier 4 (partial control)           & 0              & Multiple \\
Tier 5 (full control-flow hijack)  & 0              & 10 \\
\bottomrule
\end{tabularx}
\caption{OSS-Fuzz corpus (7\,000 entry points), reproduced
from~\citet{anthropicMythos2026}.}
\label{tab:ossfuzz}
\end{table}

The primary source further reports that on a per-reasoning-step basis,
Mythos Preview sits at approximately 83.1\,\% accuracy and Opus~4.6 at
approximately 66.6\,\%; this differential compounds geometrically over
multi-step reasoning chains~\citep{anthropicMythos2026,semgrepNeedles2026}.
At 20~steps, the implied success rates are 2.4\,\% versus 0.03\,\%,
an approximately 80$\times$ gap, which the primary source attributes to
the model checkpoint rather than to the scaffold.

\subsection{Plausibly closable gaps at the scaffold level}

We identify three categories of gap that the scaffold can plausibly
close, drawing on the public agentic-security literature.

First, the crash-finding gap (tier~1--2) is quantitative. Better file
prioritisation, more attempts via pass@$k$, VSP-focused prompting, and
tool augmentation (fuzzer, AddressSanitizer, LLM for interpretation) all
appear in reports of improved crash-finding performance.
\citet{ndss2025} show that naive zero-shot LLM prompting is close to
random on vulnerability detection and that structured approaches are
essential. A recent study of Claude Sonnet~4.5 reports a move from
12.8\,\% accuracy at the zero-day level to 95.7\,\% at full-information
level, indicating that context is the dominant factor for crash-finding.

Second, severity assessment is already close to the reported Mythos level.
The Think-\&-Verify pattern lifts accuracy from 36.7\,\% to 57.94\,\%
in~\citet{vulnSage2025}.

Third, vulnerability identification in well-documented patterns benefits
from VSP-style prompting, with \citet{vsp2024} reporting a 553\,\% F1
improvement for identification relative to generic prompting.

\subsection{Weight-level capabilities and the ceiling}

A second category of gap is, on current evidence, not closable at the
scaffold level. Multi-step exploit development (tier~3--5) requires the
model to maintain coherent reasoning state across dozens of dependent
steps: constructing a ROP chain, defeating ASLR and DEP simultaneously via
JIT heap spray, escaping a sandbox via a kernel write primitive from user
space, and performing cross-cache slab reclamation for controlled heap
state. The primary source is explicit that these capabilities emerged as
a downstream consequence of general improvements in code, reasoning, and
autonomy~\citep{anthropicGlasswing2026}. Taken literally, this places the
tier~3--5 gap in the model, not in the scaffold.

\subsection{Practical implications}

We derive several implications for practitioners. Deterministic tools
matter more than prompt engineering: Project Naptime reports a 20$\times$
improvement over the CyberSecEval2
baseline~\citep{projectzeroNaptime2024}, and this is the largest single
intervention we find in the public literature. VSP-style prompts
outperform generic ``find bugs'' prompts by a large margin. Combining
classical static analysis with LLM-based intelligent triage mirrors the
winning DARPA AIxCC configuration, where traditional fuzzing and static
analysis were augmented rather than replaced. Pass@$k$ with $k{=}20{-}50$
at \$0.05--\$0.50 per run is inexpensive insurance. A sceptical
secondary validator is a small addition with large false-positive
reduction. Ranking files before hunting avoids wasted runs. Two-pass
Think-\&-Verify reduces false positives further. On the exploit-writing
axis, we recommend recognising the weight-level ceiling and using
general-purpose models for discovery and triage while leaving exploit
development to human experts or to the next generation of specialised
checkpoints.

% ===========================================================================
\section{Ethical Considerations}
\label{sec:ethics}

The scaffold is released as a finding-first tool for self-scan and
coordinated disclosure. Four considerations are worth stating explicitly.

\textit{Scope of the public release.} The public repository contains the
scaffold code, prompts, sink catalogues, file ranker, and validator. It
does not contain the live-exploit validation stage (Phase~5), exploit-sketch
generators, or any per-target tuning accumulated during real scans. The
split is deliberate and is documented in the repository's
\texttt{SECURITY.md}.

\textit{Operator responsibility.} The intended uses are open-source-project
self-scan, security-researcher audits with explicit authorisation, and
academic replication of the methodology. Running the pipeline against
systems the operator is not authorised to audit is outside the intended
use and is the operator's responsibility.

\textit{Findings handling.} Findings surfaced by the pipeline are reported
via the upstream project's preferred security-contact channel
(GitHub Security Advisory, NVD, project \texttt{SECURITY.md}).
Proof-of-concept code is not published before the vendor has acknowledged
and patched.

\textit{Disclosure track record.} A track record of coordinated disclosures
will be maintained in the repository's \texttt{disclosures/} directory and
will be updated as embargoes lift.

% ===========================================================================
\section{Limitations and Future Work}
\label{sec:limitations}

Sink catalogues are hand-maintained and regex-based; semantic bugs
outside the catalogue are systematically missed, and the coverage of
frameworks such as GLib/GObject and \texttt{json-glib} is incomplete.
The ranker skips files over 2\,500~LOC, which is a real gap for large
parser files. The public scaffold does not currently isolate agent
execution in a sandbox; running against untrusted codebases locally is at
the operator's own risk. Most substantively, this report contains
empirical \emph{deployment} observations and cost telemetry
(\S\ref{sec:eval-cost}, \S\ref{sec:eval-yield}) but no controlled
head-to-head benchmark; the methodological reasons for that choice
are set out in \S\ref{sec:eval-why-no-bench}. Constructing an
agentic-pipeline-appropriate, training-contamination-aware benchmark
remains future work.

% ===========================================================================
\section{Conclusion}
\label{sec:conclusion}

We have described Mythos Research Edition, an open-source outside-in
replication of the scaffold described in Anthropic's April 2026 Mythos
Preview announcement, built around a general-purpose frontier model. We
have separated the scaffold contribution from the model contribution,
surveyed prior agentic security agents, delineated plausibly closable
gaps from weight-level ceilings on current evidence, and reported
aggregate cost telemetry from real-world coordinated-disclosure
deployments. We have argued in \S\ref{sec:eval-why-no-bench} that a
controlled head-to-head benchmark of agentic pipelines is itself an open
methodological problem, and that until such a benchmark exists the most
defensible empirical claims are cost telemetry, qualitative
validator-agreement observations, and coverage of finding categories.
The Research Edition is the public, finding-first subset of the
operator's workflow; downstream verification work for published
advisories uses an internal toolchain that is not part of this
artefact and is not documented here, in line with the
responsible-disclosure boundary discussed in
\S\ref{sec:approach-phase5}. The artefact, the prompts, the sink
catalogues, and the LaTeX source of this report are released under
Apache~2.0 and are archived at Zenodo for citable reference.

% ===========================================================================
\bibliographystyle{plainnat}
\bibliography{references}

\end{document}
